% Level 2, Week 9 -- Event-Driven Games (answer key)
\answersweek{9}{Event-Driven Games}{Answers for workbook chapter 9 and homework 9.}

\section{Workbook exercises}
\begin{answer}[Exercise 9.1 Etch-a-Sketch]
\textbf{What to look for:} \texttt{screen.listen()} before \texttt{turtle.done()}; every \texttt{onkey} line has the function name \emph{without} parentheses;
the key names are spelled exactly (\texttt{"Up"}, \texttt{"space"}). Any drawing description is fine.
\end{answer}
\pyfile{../exercises/week09/ex-etch-a-sketch/solution.py}

\begin{answer}[Exercise 9.2 Dodge the Meteors]
\textbf{Step 5:} logic functions are \texttt{clamp()}, \texttt{move\_player()} and \texttt{is\_touching()}: they only take numbers and return a value, so they
can be tested with pytest (like in week 8) without opening a window. Drawing functions use turtles (\texttt{make\_meteor()}, \texttt{show\_score()},
\texttt{game\_loop()}\ldots). Keeping them apart makes the rules easy to test and easy to change. Any best score is fine.
\end{answer}
\pyfile{../exercises/week09/ex-dodge-meteors/solution.py}

\begin{answer}[Reflection]
Without \texttt{ontimer} at the end of \texttt{game\_loop()}, the loop runs only once, so the meteors move one step and freeze.
\texttt{onkey(go\_left(), "Left")} calls \texttt{go\_left} straight away and gives \texttt{onkey} its answer (\texttt{None}) instead of the function.
\end{answer}

\section{Homework}
\begin{answer}[Core Part 1: Stamp Studio]
\texttt{next\_index(2, 4)} = 3 and \texttt{next\_index(3, 4)} = 0. A student who writes
\texttt{if index + 1 == length: return 0} / \texttt{return index + 1} is also correct.
\end{answer}
\pyfile{../exercises/week09/hw-stamp-studio/solution.py}

\begin{answer}[Core Part 2: Project pitch (M1)]
Good pitch pages are specific: a name, one clear sentence (\enquote{Catch falling stars with a basket and dodge the meteors}), the keys, how it ends,
and a sensible plan for the list (falling things, colors), dictionary (settings, a design), file (high scores, saved designs) and loop.
If the idea is very big, help the student find the \enquote{smallest version that is still fun} for week 10. The full M1 checklist is in the Project Teacher Guide.
\end{answer}

\begin{answer}[Challenge and Creative]
Planning page 2: look for a loop drawn as an arrow going back, at least one decision (diamond or \enquote{if}), and the end condition.
Bigger/smaller stamps: \texttt{studio["size"]} goes up or down (keep it between 1 and 5 with \texttt{clamp()}), then \texttt{stamper.shapesize(studio["size"])}.
Screen sketch: every part labeled, and parts that differ by shape or size as well as color.
\end{answer}
